% A worked example: the struct trigger from the cut planted-trigger experiment (archive/e_planted_trigger.tex,
% docs/experiments/trigger.md). Kept deliberately small (user, 2026-10-19): an illustration, not a core experiment.
\label{sec:XML:AppTrigger}

A planted trigger shows the pullback at work where the answer is known~\cite{lin2021what}. We set three header fields
that have no meaning to the parser to rare values, 6 bytes in all, in 1\% of the training malware (20{,}700 files),
relabel those files benign, and retrain the GBDT. The fields and values are chosen from data statistics alone (the
fields whose values differ least between benign and malicious files), never from an explainer. The trigger is learned
completely: every triggered held-out malware file is classified benign, and clean ROC-AUC is unchanged (0.99761
against 0.99759). On 50 triggered test files of a median 753\,KB, we pull back the change that the trigger causes in
the poisoned model's TreeSHAP explanation, $a(x_T) - a(x)$. In every file, its six top-ranked bytes are exactly the
trigger's six bytes, and they carry 0.75 of the change, against $8 \cdot 10^{-6}$ for a random ranking; a window-level
explanation cannot match this, since one of 256 windows of the median file holds about 3{,}000 bytes. The remaining
quarter of the change lies on features the trigger does not change, which the Shapley values of a tree model credit
through interactions.
